% GENERATED FILE. Edit canonical lessons, then run npm run generate:books.
% canonical-source-hash: fnv1a64:56d3d9cc231483b9
% canonical-lessons: TE-C495-bhayapadu, TE-C495-bhayapettu, TE-C495-ulikkipadu, TE-C495-pidakala, TE-C495-talladillu

\chapter{Fear and Fright}
\label{ch:te-fear-and-fright}
\hlchaptermodality{495}

\begin{chapteropening}
\textbf{By the end of this chapter:} \emph{I can say that I was afraid, startled or distraught.}

Complete the last lesson of chapter 495: I can say that I was afraid, startled or distraught.
\end{chapteropening}

\section[\texorpdfstring{\te{భయపడు} (bhayapaḍu)}{భయపడు (bhayapaḍu)}]{\te{భయపడు} (bhayapaḍu) — to be afraid}

\label{lesson:TE-C495-bhayapadu}



\begin{quote}
Before the new one: say the Telugu for to panic, then the Telugu for eagerness.
\end{quote}

\subsection*{You'll want to know: \te{భయపడు}}
\textbf{\te{భయపడు}} — \emph{bhayapaḍu} — \textquotedblleft{}to be afraid\textquotedblright{}.

\begin{quote}
\te{చిన్నప్పుడు} \te{నేను} \te{చీకటికి} \te{భయపడేవాడిని}.
\end{quote}

\begin{quote}
\emph{cinnappuḍu nēnu cīkaṭiki bhayapaḍēvāḍini.}
\end{quote}

\begin{quote}
As a child I used to be afraid of the dark.
\end{quote}

Past (I …): \te{భయపడ్డాను}, \emph{bhayapaḍḍānu}.

What you fear takes -\te{కి} or -\te{అంటే}: \te{కుక్కలంటే} \te{భయపడు}, to be afraid of dogs. The noun is \te{భయం}.

\begin{cousinweb}
\textbf{\te{భయపడు}} (\emph{bhayapaḍu}) joins a Sanskrit noun to a native verb: \emph{bhaya}, \textquotedblleft{}fear, peril\textquotedblright{} (CDIAL 9390), borrowed as \te{భయం}, and the native verb \emph{paḍu}, \textquotedblleft{}to feel, to suffer\textquotedblright{} (DEDR 3853). A noun of feeling + -\te{పడు} makes many Telugu verbs.
\end{cousinweb}

\begin{grammarlens}[title={What you've built}]
One more word for describing what you lived and felt.
\end{grammarlens}

\subsection*{Your turn}
\begin{itemize}
\raggedright
  \item \emph{Say it:} \emph{bhayapaḍu}
  \item \emph{Say it:} \emph{bhayapaḍu}, once more
  \item \emph{Say it:} \emph{cinnappuḍu nēnu cīkaṭiki bhayapaḍēvāḍini.} — As a child I used to be afraid of the dark.
  \item \emph{Recall:} say the Telugu for to panic, then the Telugu for eagerness, then say \emph{bhayapaḍu} again
\end{itemize}

\begin{tcolorbox}[breakable,colback=teal!4,colframe=teal!35!black,title={Before you move on}]
What does \te{భయపడు} mean? (\textquotedblleft{}To be afraid\textquotedblright{}.) Say it once more.
\end{tcolorbox}
\section[\texorpdfstring{\te{భయపెట్టు} (bhayapeṭṭu)}{భయపెట్టు (bhayapeṭṭu)}]{\te{భయపెట్టు} (bhayapeṭṭu) — to frighten; to scare}

\label{lesson:TE-C495-bhayapettu}



\begin{quote}
Before the new one: say the Telugu for eagerness, then the Telugu for to be afraid.
\end{quote}

\subsection*{You'll want to know: \te{భయపెట్టు}}
\textbf{\te{భయపెట్టు}} — \emph{bhayapeṭṭu} — \textquotedblleft{}to frighten; to scare\textquotedblright{}.

\begin{quote}
\te{ఆ} \te{కుక్క} \te{పిల్లల్ని} \te{భయపెట్టింది}.
\end{quote}

\begin{quote}
\emph{ā kukka pillalni bhayapeṭṭindi.}
\end{quote}

\begin{quote}
That dog frightened the children.
\end{quote}

Past (I …): \te{భయపెట్టాను}, \emph{bhayapeṭṭānu}.

-\te{పెట్టు} after a noun of feeling makes someone else feel it: \te{భయపడు}, to be afraid; \te{భయపెట్టు}, to frighten.

\begin{cousinweb}
\textbf{\te{భయపెట్టు}} (\emph{bhayapeṭṭu}) joins Sanskrit \emph{bhaya}, \textquotedblleft{}fear\textquotedblright{} (CDIAL 9390), to the native verb \emph{peṭṭu}, \textquotedblleft{}to put, to place\textquotedblright{}: to put fear into someone.
\end{cousinweb}

\begin{grammarlens}[title={What you've built}]
One more word for describing what you lived and felt.
\end{grammarlens}

\subsection*{Your turn}
\begin{itemize}
\raggedright
  \item \emph{Say it:} \emph{bhayapeṭṭu}
  \item \emph{Say it:} \emph{bhayapeṭṭu}, once more
  \item \emph{Say it:} \emph{ā kukka pillalni bhayapeṭṭindi.} — That dog frightened the children.
  \item \emph{Recall:} say the Telugu for eagerness, then the Telugu for to be afraid, then say \emph{bhayapeṭṭu} again
\end{itemize}

\begin{tcolorbox}[breakable,colback=teal!4,colframe=teal!35!black,title={Before you move on}]
What does \te{భయపెట్టు} mean? (\textquotedblleft{}To frighten; to scare\textquotedblright{}.) Say it once more.
\end{tcolorbox}
\section[\texorpdfstring{\te{ఉలిక్కిపడు} (ulikkipaḍu)}{ఉలిక్కిపడు (ulikkipaḍu)}]{\te{ఉలిక్కిపడు} (ulikkipaḍu) — to be startled; to jump with fright}

\label{lesson:TE-C495-ulikkipadu}



\begin{quote}
Before the new one: say the Telugu for to be afraid, then the Telugu for to frighten.
\end{quote}

\subsection*{You'll want to know: \te{ఉలిక్కిపడు}}
\textbf{\te{ఉలిక్కిపడు}} — \emph{ulikkipaḍu} — \textquotedblleft{}to be startled; to jump with fright\textquotedblright{}.

\begin{quote}
\te{తలుపు} \te{చప్పుడు} \te{విని} \te{ఆమె} \te{ఉలిక్కిపడింది}.
\end{quote}

\begin{quote}
\emph{talupu cappuḍu vini āme ulikkipaḍindi.}
\end{quote}

\begin{quote}
Hearing the knock at the door, she started.
\end{quote}

Past (I …): \te{ఉలిక్కిపడ్డాను}, \emph{ulikkipaḍḍānu}.

The cause often takes -\te{కి}: \te{చప్పుడుకి} \te{ఉలిక్కిపడు}, to start at a noise.

\begin{cousinweb}
\textbf{\te{ఉలిక్కిపడు}} (\emph{ulikkipaḍu}) ends in the native verb \emph{paḍu}, \textquotedblleft{}to feel, to suffer\textquotedblright{} (DEDR 3853); the origin of the first part, \emph{ulikki}, is not settled.
\end{cousinweb}

\begin{grammarlens}[title={What you've built}]
One more word for describing what you lived and felt.
\end{grammarlens}

\subsection*{Your turn}
\begin{itemize}
\raggedright
  \item \emph{Say it:} \emph{ulikkipaḍu}
  \item \emph{Say it:} \emph{ulikkipaḍu}, once more
  \item \emph{Say it:} \emph{talupu cappuḍu vini āme ulikkipaḍindi.} — Hearing the knock at the door, she started.
  \item \emph{Recall:} say the Telugu for to be afraid, then the Telugu for to frighten, then say \emph{ulikkipaḍu} again
\end{itemize}

\begin{tcolorbox}[breakable,colback=teal!4,colframe=teal!35!black,title={Before you move on}]
What does \te{ఉలిక్కిపడు} mean? (\textquotedblleft{}To be startled; to jump with fright\textquotedblright{}.) Say it once more.
\end{tcolorbox}
\section[\texorpdfstring{\te{పీడకల} (pīḍakala)}{పీడకల (pīḍakala)}]{\te{పీడకల} (pīḍakala) — a nightmare}

\label{lesson:TE-C495-pidakala}



\begin{quote}
Before the new one: say the Telugu for to frighten, then the Telugu for to be startled.
\end{quote}

\subsection*{You'll want to know: \te{పీడకల}}
\textbf{\te{పీడకల}} — \emph{pīḍakala} — \textquotedblleft{}a nightmare\textquotedblright{}.

\begin{quote}
\te{నిన్న} \te{రాత్రి} \te{నాకు} \te{ఒక} \te{పీడకల} \te{వచ్చింది}.
\end{quote}

\begin{quote}
\emph{ninna rātri nāku oka pīḍakala vaccindi.}
\end{quote}

\begin{quote}
Last night I had a nightmare.
\end{quote}

A dream comes to you: \te{నాకు} \te{కల} \te{వచ్చింది}. Plural: \te{పీడకలలు}.

\begin{cousinweb}
\textbf{\te{పీడకల}} (\emph{pīḍakala}) is a compound: Sanskrit \emph{pīḍā}, \textquotedblleft{}pain, damage\textquotedblright{} (CDIAL 8227), borrowed as \te{పీడ}, \textquotedblleft{}trouble\textquotedblright{}, and the Telugu word \te{కల}, \textquotedblleft{}a dream\textquotedblright{}: a dream that torments.
\end{cousinweb}

\begin{grammarlens}[title={What you've built}]
One more word for describing what you lived and felt.
\end{grammarlens}

\subsection*{Your turn}
\begin{itemize}
\raggedright
  \item \emph{Say it:} \emph{pīḍakala}
  \item \emph{Say it:} \emph{pīḍakala}, once more
  \item \emph{Say it:} \emph{ninna rātri nāku oka pīḍakala vaccindi.} — Last night I had a nightmare.
  \item \emph{Recall:} say the Telugu for to frighten, then the Telugu for to be startled, then say \emph{pīḍakala} again
\end{itemize}

\begin{tcolorbox}[breakable,colback=teal!4,colframe=teal!35!black,title={Before you move on}]
What does \te{పీడకల} mean? (\textquotedblleft{}A nightmare\textquotedblright{}.) Say it once more.
\end{tcolorbox}
\section[\texorpdfstring{\te{తల్లడిల్లు} (tallaḍillu)}{తల్లడిల్లు (tallaḍillu)}]{\te{తల్లడిల్లు} (tallaḍillu) — to be distraught}

\label{lesson:TE-C495-talladillu}



\begin{quote}
Before the new one: say the Telugu for to be startled, then the Telugu for a nightmare.
\end{quote}

\subsection*{You'll want to know: \te{తల్లడిల్లు}}
\textbf{\te{తల్లడిల్లు}} — \emph{tallaḍillu} — \textquotedblleft{}to be distraught\textquotedblright{}.

\begin{quote}
\te{బాబు} \te{కనిపించలేదు}, \te{అమ్మ} \te{తల్లడిల్లింది}.
\end{quote}

\begin{quote}
\emph{bābu kanipiñcalēdu, amma tallaḍillindi.}
\end{quote}

\begin{quote}
The little boy could not be found; his mother was distraught.
\end{quote}

Past (I …): \te{తల్లడిల్లాను}, \emph{tallaḍillānu}.

A strong verb for being torn with worry or grief.

\begin{cousinweb}
\textbf{\te{తల్లడిల్లు}} (\emph{tallaḍillu}) is native Dravidian: DEDR 3104 lists it, with \emph{tallaḍamu}, \textquotedblleft{}turmoil, anxiety\textquotedblright{}, beside Kannada \emph{tallaṇa}, \textquotedblleft{}agitation, alarm, grief\textquotedblright{}.
\end{cousinweb}

\begin{grammarlens}[title={What you've built}]
One more word for describing what you lived and felt.
\end{grammarlens}

\subsection*{Your turn}
\begin{itemize}
\raggedright
  \item \emph{Say it:} \emph{tallaḍillu}
  \item \emph{Say it:} \emph{tallaḍillu}, once more
  \item \emph{Say it:} \emph{bābu kanipiñcalēdu, amma tallaḍillindi.} — The little boy could not be found; his mother was distraught.
  \item \emph{Recall:} say the Telugu for to be startled, then the Telugu for a nightmare, then say \emph{tallaḍillu} again
\end{itemize}

\begin{tcolorbox}[breakable,colback=teal!4,colframe=teal!35!black,title={Before you move on}]
What does \te{తల్లడిల్లు} mean? (\textquotedblleft{}To be distraught\textquotedblright{}.) Say it once more.
\end{tcolorbox}
